\section{Results to be completed for the 22-family model}
\label{sec:results}
\TwentyTwoScope

\subsection{Evaluation population and exact assembly}
\ResultSlot{R01--R03}{Insert the admitted cohort and frozen held-out evaluation population, final checkpoint and inference-policy identifiers, attempt budget, independent training seeds and uncertainty convention. No development result is a substitute for this experiment.}
All requested outputs will contribute to exact-L1 yield, including invalid graphs, failed searches,
source-domain rejections and abstentions. Report validity and connectedness per request, decomposition coverage,
candidate precision and ambiguity separately. Equal-family averages and pooled counts accompany
all individual families, and cannot conceal a weak family. The final seed summaries will use means
and sample standard deviations across independent training runs; the planned minimum is three.
Sampling variation from a fixed checkpoint will be reported separately. Tables in
Appendix~\ref{app:22-results} specify the remaining endpoints and denominators.
\begin{table}[H]
\caption{\textbf{Exact L1 across all 22 assembly families, final results pending.} All cells require independently trained seed summaries under matched frozen requests. ``Final'' includes the declared checked inference procedure; raw and final outputs share attempt IDs. Null and cyclic controls require their own matched training. Candidate precision is distinct from replay of already accepted traces. Per-seed counts and denominator definitions: Appendix~\ref{app:22-results}.}
\label{tab:22-l1}
\centering\footnotesize\setlength{\tabcolsep}{3pt}\renewcommand{\arraystretch}{1.12}
\begin{tabular}{@{}>{\raggedright\arraybackslash}p{.30\textwidth}ccccc@{}}
\toprule
Family & \shortstack{Raw\\L1 (\%)} & \shortstack{Final\\L1 (\%)} & \shortstack{Shared\\null (\%)} & \shortstack{Cyclic\\label (\%)} & \shortstack{Decomp. cov. /\\candidate prec.} \\
\midrule
A3 amine--aldehyde--alkyne & \PendingCell & \PendingCell & \PendingCell & \PendingCell & \PendingCell \\
Acid--epoxide diester & \PendingCell & \PendingCell & \PendingCell & \PendingCell & \PendingCell \\
AEMA & \PendingCell & \PendingCell & \PendingCell & \PendingCell & \PendingCell \\
Aldehyde Ugi 3-CR & \PendingCell & \PendingCell & \PendingCell & \PendingCell & \PendingCell \\
Aldehyde Ugi 4-CR & \PendingCell & \PendingCell & \PendingCell & \PendingCell & \PendingCell \\
$\alpha$-Isocyanoester dihydroimidazole & \PendingCell & \PendingCell & \PendingCell & \PendingCell & \PendingCell \\
Amine alkylation & \PendingCell & \PendingCell & \PendingCell & \PendingCell & \PendingCell \\
Amine--epoxide & \PendingCell & \PendingCell & \PendingCell & \PendingCell & \PendingCell \\
Aryl reductive amination & \PendingCell & \PendingCell & \PendingCell & \PendingCell & \PendingCell \\
Aza-Michael acrylamide & \PendingCell & \PendingCell & \PendingCell & \PendingCell & \PendingCell \\
Aza-Michael acrylate & \PendingCell & \PendingCell & \PendingCell & \PendingCell & \PendingCell \\
Disulfide Michael & \PendingCell & \PendingCell & \PendingCell & \PendingCell & \PendingCell \\
Epoxide O-acylation & \PendingCell & \PendingCell & \PendingCell & \PendingCell & \PendingCell \\
iPhos & \PendingCell & \PendingCell & \PendingCell & \PendingCell & \PendingCell \\
Ketone--isocyanide amide & \PendingCell & \PendingCell & \PendingCell & \PendingCell & \PendingCell \\
Ketone Ugi 4-CR & \PendingCell & \PendingCell & \PendingCell & \PendingCell & \PendingCell \\
Maleate & \PendingCell & \PendingCell & \PendingCell & \PendingCell & \PendingCell \\
O-esterification & \PendingCell & \PendingCell & \PendingCell & \PendingCell & \PendingCell \\
Passerini & \PendingCell & \PendingCell & \PendingCell & \PendingCell & \PendingCell \\
Preassembled thiol-yne & \PendingCell & \PendingCell & \PendingCell & \PendingCell & \PendingCell \\
Reductive amination & \PendingCell & \PendingCell & \PendingCell & \PendingCell & \PendingCell \\
Thiolactone & \PendingCell & \PendingCell & \PendingCell & \PendingCell & \PendingCell \\
\midrule
Equal-family macro & \PendingCell & \PendingCell & \PendingCell & \PendingCell & \PendingCell \\
Pooled all requests & \PendingCell & \PendingCell & \PendingCell & \PendingCell & \PendingCell \\
Worst-family result & \PendingCell & \PendingCell & \PendingCell & \PendingCell & \PendingCell \\
\bottomrule
\end{tabular}
\end{table}

\subsection{Structural quality, realism and component diversity}
\ResultSlot{R06--R09}{Insert family-wise representation failures, chemical alerts and context flags, qualified head/ring agreement, reference distribution comparisons, component diversity and held-component recovery. State residual failures and abstentions alongside improvements.}
Exact L1 is a reaction-consistency result. Structural assessment separately checks the complete final
graph, including retained headgroup sites, hydrophobic carbon-domain organization and ring-system
connectivity. A basic-substructure match does not establish ionization. Absence of a narrow chemical
alert does not establish high overall quality. Comparison with TRAIN local features is descriptive
when those same features supply a repair prior; independent source-held-out controls provide a
separate assessment. Product and role-specific component diversity and novelty are retained as
separate outcomes, rather than combined into a single quality score.

\begin{figure}[H]
\centering
\DraftPanel{a. Assembly and conditions}{Per-family raw/final exact yield, request counts and training-seed intervals.}\hfill
\DraftPanel{b. Structural assessment}{Chemical alerts, qualified condition failures and explicit abstentions.}\\[6pt]
\DraftPanel{c. Diversity and realism}{Role diversity, novelty and independent reference precision/coverage.}\hfill
\DraftPanel{d. Evidence depth}{L1, L2 and L3 dispositions; complete dossier yield per request.}
\caption{\textbf{Planned 22-family evidence overview.} These empty panels reserve space for final
measurements. Each will include all families, fixed comparator budgets and its own denominator.
No illustrative trend, fabricated value or implied positive outcome is shown.}
\label{fig:22-overview}
\end{figure}

\subsection{Attribution, efficiency and synthesis evidence}
\ResultSlot{R04--R05, R10--R12}{Insert matched conditioning controls, the original-three-family bridge,
decoder and architecture ablations, training adequacy, cost and complete-route results. Distinguish
same-cost contrasts from improvements purchased by additional proposals or source checks.}
A learned-conditioning claim requires matched independently trained controls and an inventory of
reaction information supplied through layouts, masks, cores, adapters and decoder priors. Fixed-weight
repairs isolate an inference contribution; they do not establish improved learning. The route analysis
will report exact precursor preparation evidence (L2), qualified terminal identity and availability
(L3), and complete-product dossiers. An upstream step or a verified assembly alone does not close
the complete route. The manuscript will retain unresolved and search-censored outcomes, including
families whose evidence is shallower than the AGILE-type Ugi case.
\ResultSlot{R13--R15}{Add a seeded all-family structure atlas, corrected and unresolved failure cases,
the final implementation/theory correspondence, and any authorized guidance contrast whose
readiness gates pass. Guidance is conditional and no wet-lab study is proposed here.}
